\chapter{L15 -- Osservabilità, ricostruibilità e forma standard di osservabilità}

\section{Idea generale della lezione}

Dopo aver studiato la raggiungibilità, passiamo al problema duale: capire quante informazioni sullo stato si possono ricavare osservando l'uscita del sistema.

Le due domande fondamentali sono:

\begin{itemize}
    \item \textbf{Osservabilità}: conoscendo \(u(\cdot)\) e \(y(\cdot)\), posso ricostruire lo \emph{stato iniziale} \(x(0)\)?
    \item \textbf{Ricostruibilità}: conoscendo \(u(\cdot)\) e \(y(\cdot)\), posso almeno ricostruire lo \emph{stato attuale} \(x(t)\), anche se lo stato iniziale non è identificabile in modo univoco?
\end{itemize}

Questa distinzione è importante soprattutto in tempo discreto.  
In tempo continuo, come vedremo, osservabilità e ricostruibilità finiscono per coincidere.

\bigskip

La logica è perfettamente duale rispetto alla raggiungibilità:

\[
\text{raggiungibilità} \quad \Longleftrightarrow \quad \text{posso muovere lo stato con l'ingresso},
\]
\[
\text{osservabilità} \quad \Longleftrightarrow \quad \text{posso ricostruire lo stato guardando l'uscita}.
\]

%=========================================================
\section{Osservabilità in tempo continuo: richiamo}

Consideriamo un sistema lineare tempo invariante in tempo continuo:
\[
\dot x(t)=Ax(t)+Bu(t),\qquad y(t)=Cx(t)+Du(t).
\]

Fissiamo due stati iniziali \(x_0\) e \(\bar x_0\).  
Diremo che \(x_0\) e \(\bar x_0\) sono \textbf{indistinguibili} in un intervallo \([0,t_f]\) se, a parità di ingresso, producono la stessa uscita per tutto l'intervallo:
\[
y(\tau; x_0,u)=y(\tau; \bar x_0,u)
\qquad
\forall \tau\in[0,t_f],\ \forall u.
\]

Poiché la soluzione dell'uscita è
\[
y(\tau)=Ce^{A\tau}x_0+\int_0^\tau Ce^{A(\tau-\sigma)}Bu(\sigma)\,d\sigma + Du(\tau),
\]
sottraendo le due uscite si ottiene
\[
Ce^{A\tau}(x_0-\bar x_0)=0
\qquad
\forall \tau\in[0,t_f].
\]

Quindi il sottospazio degli stati indistinguibili dall'origine è
\[
\mathcal I_{t_f}(0)
=
\bigcap_{\tau\in[0,t_f]} \ker\!\big(Ce^{A\tau}\big).
\]

Dalla teoria già vista si sa che questo insieme coincide con il nucleo della matrice di osservabilità classica:
\[
\mathcal O=
\begin{bmatrix}
C\\
CA\\
CA^2\\
\vdots\\
CA^{n-1}
\end{bmatrix}.
\]

In particolare:
\[
\mathcal I_{t_f}(0)=\ker(\mathcal O).
\]

Quindi il sistema è \textbf{completamente osservabile} se e solo se
\[
\ker(\mathcal O)=\{0\}
\qquad
\Longleftrightarrow
\qquad
\rank(\mathcal O)=n.
\]

\medskip

Osservazione importante: nel caso continuo, il concetto di osservabilità è una proprietà strutturale del sistema e non dipende dalla scelta dell'istante finale \(t_f>0\).

%=========================================================
\section{Osservabilità in tempo discreto}

Consideriamo ora un sistema lineare tempo invariante in tempo discreto:
\[
\begin{cases}
x(k+1)=Ax(k)+Bu(k),\\[2mm]
y(k)=Cx(k)+Du(k).
\end{cases}
\]

La risposta del sistema al tempo \(\tau\), partendo da \(x_0\) all'istante \(0\), è:
\[
y(\tau;0,x_0,u)
=
CA^\tau x_0
+
\sum_{k=0}^{\tau-1} CA^{\tau-k-1}Bu(k)
+
Du(\tau),
\qquad \tau\in\{0,1,2,\dots\}.
\]

Se due stati iniziali \(x_0\) e \(\bar x_0\) sono indistinguibili in \(t\) passi, allora per ogni \(\tau=0,\dots,t-1\) deve valere
\[
y(\tau;0,x_0,u)=y(\tau;0,\bar x_0,u).
\]

Sottraendo membro a membro si cancellano tutti i termini forzati e rimane:
\[
\Delta y(\tau)
=
y(\tau;0,x_0,u)-y(\tau;0,\bar x_0,u)
=
CA^\tau (x_0-\bar x_0).
\]

Se ora raccogliamo tutte le differenze di uscita in un unico vettore, otteniamo
\[
\begin{bmatrix}
\Delta y(0)\\
\Delta y(1)\\
\Delta y(2)\\
\vdots\\
\Delta y(t-1)
\end{bmatrix}
=
\underbrace{
\begin{bmatrix}
C\\
CA\\
CA^2\\
\vdots\\
CA^{t-1}
\end{bmatrix}}_{\displaystyle \mathcal O_t}
(x_0-\bar x_0).
\]

La matrice
\[
\mathcal O_t=
\begin{bmatrix}
C\\
CA\\
CA^2\\
\vdots\\
CA^{t-1}
\end{bmatrix}
\]
si chiama \textbf{matrice di osservabilità in \(t\) osservazioni}.

\subsection{Sottospazio di inosservabilità in \(t\) passi}

Definiamo
\[
\bar{\mathcal O}_t:=\ker(\mathcal O_t).
\]
Questo è il \textbf{sottospazio di inosservabilità in \(t\) passi}.

Esso contiene tutte le differenze iniziali \(x_0-\bar x_0\) che, per \(t\) passi, non lasciano alcuna traccia in uscita.

In altri termini:
\[
x_0 \sim_t \bar x_0
\qquad \Longleftrightarrow \qquad
x_0-\bar x_0\in \bar{\mathcal O}_t.
\]

Quindi:
\[
\bar{\mathcal O}_t=\{0\}
\qquad \Longleftrightarrow \qquad
\text{il sistema è completamente osservabile in \(t\) passi.}
\]

%=========================================================
\section{Catena dei sottospazi di inosservabilità}

Aumentando il numero di osservazioni, si aggiungono righe alla matrice \(\mathcal O_t\).  
Di conseguenza il nucleo può solo restringersi:
\[
\bar{\mathcal O}_1 \supseteq \bar{\mathcal O}_2 \supseteq \bar{\mathcal O}_3 \supseteq \cdots
\]

Questo è intuitivo: più a lungo osservo il sistema, più informazione raccolgo, quindi meno stati iniziali restano indistinguibili.

\subsection{Quando si ferma questa catena?}

La risposta è: \emph{al più dopo \(n\) osservazioni}.

Infatti, per il teorema di Cayley--Hamilton, ogni potenza \(A^k\) con \(k\ge n\) si può scrivere come combinazione lineare di
\[
I,\ A,\ A^2,\ \dots,\ A^{n-1}.
\]
Quindi anche ogni riga \(CA^k\) con \(k\ge n\) è combinazione lineare delle righe precedenti.  
Ne segue che oltre \(n\) osservazioni non compare più nessuna informazione nuova.

Perciò:
\[
\bar{\mathcal O}_n=\bar{\mathcal O}_{n+1}=\bar{\mathcal O}_{n+2}=\cdots
\]

e dunque, per verificare la completa osservabilità, basta controllare
\[
\rank(\mathcal O_n)=n.
\]

%=========================================================
\section{Primo esempio: nessun miglioramento aumentando le osservazioni}

Consideriamo
\[
A=
\begin{bmatrix}
-1 & 2 & 0\\
0 & -1 & 0\\
1 & -1 & 2
\end{bmatrix},
\qquad
C=
\begin{bmatrix}
0 & 1 & 0
\end{bmatrix}.
\]

\subsection{Prima osservazione}

Si ha
\[
\mathcal O_1=C=
\begin{bmatrix}
0 & 1 & 0
\end{bmatrix}.
\]
Quindi
\[
\bar{\mathcal O}_1=\ker(\mathcal O_1)
=
\left\{
\begin{bmatrix}
\alpha\\
0\\
\beta
\end{bmatrix}
:\ \alpha,\beta\in\mathbb R
\right\}
=
\operatorname{span}\left\{
\begin{bmatrix}
1\\0\\0
\end{bmatrix},
\begin{bmatrix}
0\\0\\1
\end{bmatrix}
\right\}.
\]
Dopo un solo passo, quindi, conosco solo la seconda componente dello stato.

\subsection{Seconda osservazione}

Calcoliamo
\[
CA=
\begin{bmatrix}
0 & 1 & 0
\end{bmatrix}
\begin{bmatrix}
-1 & 2 & 0\\
0 & -1 & 0\\
1 & -1 & 2
\end{bmatrix}
=
\begin{bmatrix}
0 & -1 & 0
\end{bmatrix}.
\]
Quindi
\[
\mathcal O_2=
\begin{bmatrix}
C\\
CA
\end{bmatrix}
=
\begin{bmatrix}
0 & 1 & 0\\
0 & -1 & 0
\end{bmatrix}.
\]
Ma le due righe sono linearmente dipendenti, quindi
\[
\bar{\mathcal O}_2=\bar{\mathcal O}_1.
\]

\subsection{Terza osservazione}

Analogamente:
\[
CA^2=
\begin{bmatrix}
0 & 1 & 0
\end{bmatrix},
\]
quindi anche
\[
\bar{\mathcal O}_3=\bar{\mathcal O}_2.
\]

\medskip

Conclusione: in questo esempio
\[
\bar{\mathcal O}_1=\bar{\mathcal O}_2=\bar{\mathcal O}_3,
\]
cioè già dopo la prima osservazione ho raccolto tutta l'informazione che l'uscita può dare.

%=========================================================
\section{Secondo esempio: le osservazioni riducono davvero il sottospazio di inosservabilità}

Consideriamo lo stesso \(A\), ma adesso
\[
C=
\begin{bmatrix}
1 & 1 & 0
\end{bmatrix}.
\]

\subsection{Prima osservazione}

La matrice di osservabilità a un passo è
\[
\mathcal O_1=
\begin{bmatrix}
1 & 1 & 0
\end{bmatrix}.
\]
Quindi
\[
\bar{\mathcal O}_1=\ker(\mathcal O_1)
=
\left\{
\begin{bmatrix}
\alpha\\
-\alpha\\
\beta
\end{bmatrix}
:\ \alpha,\beta\in\mathbb R
\right\}
=
\operatorname{span}\left\{
\begin{bmatrix}
1\\-1\\0
\end{bmatrix},
\begin{bmatrix}
0\\0\\1
\end{bmatrix}
\right\}.
\]

\subsection{Seconda osservazione}

Calcoliamo
\[
CA=
\begin{bmatrix}
1 & 1 & 0
\end{bmatrix}
\begin{bmatrix}
-1 & 2 & 0\\
0 & -1 & 0\\
1 & -1 & 2
\end{bmatrix}
=
\begin{bmatrix}
-1 & 1 & 0
\end{bmatrix}.
\]
Quindi
\[
\mathcal O_2=
\begin{bmatrix}
1 & 1 & 0\\
-1 & 1 & 0
\end{bmatrix}.
\]
La condizione
\[
\mathcal O_2 x=0
\]
equivale al sistema
\[
\begin{cases}
x_1+x_2=0,\\
-x_1+x_2=0,
\end{cases}
\]
da cui segue
\[
x_1=x_2=0.
\]
Quindi
\[
\bar{\mathcal O}_2=
\operatorname{span}\left\{
\begin{bmatrix}
0\\0\\1
\end{bmatrix}
\right\}.
\]

\subsection{Terza osservazione}

Calcoliamo ancora
\[
CA^2=
\begin{bmatrix}
1 & -3 & 0
\end{bmatrix}.
\]
Quindi
\[
\mathcal O_3=
\begin{bmatrix}
1 & 1 & 0\\
-1 & 1 & 0\\
1 & -3 & 0
\end{bmatrix}.
\]
Ma la terza riga non introduce una nuova direzione indipendente, quindi
\[
\bar{\mathcal O}_3=\bar{\mathcal O}_2.
\]

\medskip

Conclusione: in questo caso
\[
\dim \bar{\mathcal O}_1=2,\qquad
\dim \bar{\mathcal O}_2=1,\qquad
\dim \bar{\mathcal O}_3=1.
\]
Dunque la seconda osservazione riduce il sottospazio di inosservabilità, mentre la terza non aggiunge più informazione.

%=========================================================
\section{Ricostruibilità in tempo discreto}

L'osservabilità riguarda la ricostruzione dello \emph{passato}, cioè dello stato iniziale \(x(0)\).  
La ricostruibilità riguarda invece lo \emph{stato attuale} \(x(t)\).

\medskip

La domanda è:

\begin{quote}
Se due stati iniziali \(x_0\) e \(\bar x_0\) sono indistinguibili in \(t\) passi, allora dopo \(t\) passi conducono necessariamente allo stesso stato?
\end{quote}

Se la risposta è sì, allora anche se non so distinguere esattamente da quale stato iniziale sono partito, posso comunque sapere in modo univoco dove si trova il sistema al tempo \(t\).

\subsection{Definizione}

Il sistema è detto \textbf{ricostruibile in \(t\) passi} se per ogni coppia di stati indistinguibili \(x_0,\bar x_0\) in \(t\) passi si ha
\[
x(t;x_0,u)=x(t;\bar x_0,u)
\qquad
\forall u.
\]

\subsection{Condizione equivalente}

Le due traiettorie sono
\[
x_1(t)=A^t x_0+\sum_{k=0}^{t-1} A^{t-k-1}Bu(k),
\]
\[
x_2(t)=A^t \bar x_0+\sum_{k=0}^{t-1} A^{t-k-1}Bu(k).
\]
Sottraendo:
\[
x_2(t)-x_1(t)=A^t(\bar x_0-x_0).
\]

Se \(x_0\) e \(\bar x_0\) sono indistinguibili in \(t\) passi, allora
\[
\bar x_0-x_0\in \ker(\mathcal O_t)=\bar{\mathcal O}_t.
\]
Perché gli stati al tempo \(t\) coincidano, serve allora che
\[
A^t(\bar x_0-x_0)=0.
\]
Quindi il sistema è ricostruibile in \(t\) passi se e solo se
\[
\boxed{
\ker(\mathcal O_t)\subseteq \ker(A^t).
}
\]

Questa è la caratterizzazione fondamentale della ricostruibilità nel discreto.

\subsection{Relazione con l'osservabilità}

Se il sistema è osservabile in \(t\) passi, allora
\[
\ker(\mathcal O_t)=\{0\},
\]
e dunque automaticamente
\[
\ker(\mathcal O_t)\subseteq \ker(A^t).
\]
Quindi:
\[
\boxed{
\text{osservabilità} \ \Longrightarrow\ \text{ricostruibilità}.
}
\]

Il viceversa, in generale, \textbf{non è vero} nel tempo discreto.

\subsection{Quando coincidono?}

Se \(A\) è invertibile, allora anche \(A^t\) è invertibile e quindi
\[
\ker(A^t)=\{0\}.
\]
In questo caso la condizione
\[
\ker(\mathcal O_t)\subseteq \ker(A^t)
\]
diventa
\[
\ker(\mathcal O_t)\subseteq \{0\},
\]
cioè
\[
\ker(\mathcal O_t)=\{0\}.
\]
Quindi, se \(A\) è invertibile,
\[
\boxed{
\text{ricostruibilità} \ \Longleftrightarrow\ \text{osservabilità}.
}
\]

\subsection{Caso continuo}

Nel caso continuo, al posto di \(A^t\) compare \(e^{At}\), ma \(e^{At}\) è sempre invertibile.  
Per questo motivo, in tempo continuo, osservabilità e ricostruibilità coincidono sempre.

\[
\boxed{
\text{nel tempo continuo: ricostruibilità } \Longleftrightarrow \text{ osservabilità.}
}
\]

%=========================================================
\section{Invarianza per cambi di coordinate}

Come per la raggiungibilità, anche l'osservabilità è una proprietà strutturale del sistema e non dipende dalla scelta delle coordinate di stato.

Consideriamo il cambio di coordinate
\[
x=Tz,
\qquad T \ \text{invertibile}.
\]
Allora il sistema diventa
\[
z(k+1)=T^{-1}AT\,z(k)+T^{-1}Bu(k),
\qquad
y(k)=CT\,z(k)+Du(k).
\]

La nuova matrice di osservabilità è
\[
\mathcal O_z=
\begin{bmatrix}
CT\\
CTT^{-1}AT\\
CT(T^{-1}AT)^2\\
\vdots\\
CT(T^{-1}AT)^{n-1}
\end{bmatrix}.
\]
Semplificando:
\[
\mathcal O_z=
\begin{bmatrix}
C\\
CA\\
CA^2\\
\vdots\\
CA^{n-1}
\end{bmatrix}T
=
\mathcal O_x\,T.
\]

Quindi
\[
\rank(\mathcal O_z)=\rank(\mathcal O_x),
\]
e inoltre
\[
\ker(\mathcal O_z)=T^{-1}\ker(\mathcal O_x).
\]

Conclusione:
\[
\boxed{
\text{l'osservabilità è una proprietà strutturale del sistema.}
}
\]

%=========================================================
\section{Dualità tra raggiungibilità e osservabilità}

Questa è una delle idee più importanti del corso.

Consideriamo il sistema
\[
\begin{cases}
x^+=Ax+Bu,\\
y=Cx+Du.
\end{cases}
\]
Il \textbf{sistema duale} è
\[
\begin{cases}
\xi^+=A^\top \xi + C^\top v,\\
\eta=B^\top \xi + D^\top v.
\end{cases}
\]

La matrice di raggiungibilità del sistema duale è
\[
R^\star=
\begin{bmatrix}
C^\top & A^\top C^\top & (A^\top)^2 C^\top & \cdots & (A^\top)^{n-1}C^\top
\end{bmatrix}.
\]
Ma questa è esattamente la trasposta della matrice di osservabilità del sistema originale:
\[
(R^\star)^\top=
\begin{bmatrix}
C\\
CA\\
CA^2\\
\vdots\\
CA^{n-1}
\end{bmatrix}
=
\mathcal O.
\]

Di conseguenza:
\[
\boxed{
(A,C)\ \text{osservabile}
\quad \Longleftrightarrow \quad
(A^\top,C^\top)\ \text{raggiungibile}.
}
\]

Questo parallelismo permette di trasferire quasi automaticamente molti risultati dalla teoria della raggiungibilità a quella dell'osservabilità.

%=========================================================
\section{Teorema fondamentale: caratterizzazione del sottospazio di inosservabilità}

Sia
\[
\bar{\mathcal O}:=\ker(\mathcal O)
\]
il sottospazio di inosservabilità del sistema.

Vale il seguente risultato fondamentale:

\medskip

\noindent
\textbf{Teorema.}
\emph{
\(\bar{\mathcal O}\) è il più grande sottospazio \(A\)-invariante contenuto in \(\ker(C)\).
}

\subsection*{Dimostrazione}

La dimostrazione ha due parti.

\paragraph{1) \(\bar{\mathcal O}\) è contenuto in \(\ker(C)\) ed è \(A\)-invariante.}

Se \(x\in\bar{\mathcal O}\), allora
\[
\mathcal O x=0,
\]
cioè
\[
Cx=0,\qquad CAx=0,\qquad CA^2x=0,\qquad \dots,\qquad CA^{n-1}x=0.
\]
In particolare \(Cx=0\), quindi
\[
x\in\ker(C).
\]

Mostriamo ora l'invarianza.  
Se \(x\in\bar{\mathcal O}\), allora
\[
CA^k x=0\qquad \forall k=0,\dots,n-1.
\]
Vogliamo verificare che anche \(Ax\in\bar{\mathcal O}\), cioè che
\[
CA^k(Ax)=0
\qquad \forall k=0,\dots,n-1.
\]
Ma
\[
CA^k(Ax)=CA^{k+1}x.
\]
Per \(k=0,\dots,n-2\) questo è nullo per definizione di \(\bar{\mathcal O}\).  
Per \(k=n-1\) compare \(CA^n x\), ma per Cayley--Hamilton \(A^n\) è combinazione lineare di \(I,A,\dots,A^{n-1}\), quindi anche \(CA^n x\) è combinazione lineare di
\[
Cx,\ CAx,\ \dots,\ CA^{n-1}x,
\]
tutti nulli. Dunque
\[
CA^n x=0.
\]
Quindi \(Ax\in\bar{\mathcal O}\), ossia \(\bar{\mathcal O}\) è \(A\)-invariante.

\paragraph{2) È il più grande con questa proprietà.}

Sia \(S\subseteq \mathbb R^n\) un sottospazio tale che:
\[
S\subseteq \ker(C),
\qquad
AS\subseteq S.
\]
Prendiamo \(x\in S\). Siccome \(S\subseteq\ker(C)\), si ha \(Cx=0\).  
Poiché \(AS\subseteq S\), allora \(Ax\in S\), e quindi
\[
CAx=0.
\]
Ripetendo:
\[
A^2x\in S \Rightarrow CA^2x=0,
\qquad \dots,
\qquad A^{n-1}x\in S \Rightarrow CA^{n-1}x=0.
\]
Dunque
\[
\mathcal O x=0,
\]
cioè \(x\in \ker(\mathcal O)=\bar{\mathcal O}\).

Quindi
\[
S\subseteq \bar{\mathcal O}.
\]
Abbiamo dunque provato che \(\bar{\mathcal O}\) è il più grande sottospazio \(A\)-invariante contenuto in \(\ker(C)\). \qed

%=========================================================
\section{Forma standard di osservabilità}

Il teorema precedente suggerisce di scegliere nuove coordinate che separino esplicitamente la parte osservabile da quella non osservabile.

\subsection{Scelta della base}

Sia
\[
\bar o := \dim(\bar{\mathcal O}).
\]
Scegliamo:
\begin{itemize}
    \item \(T_{\bar o}\in\mathbb R^{n\times \bar o}\): una matrice le cui colonne formano una base di \(\bar{\mathcal O}\);
    \item \(T_o\in\mathbb R^{n\times (n-\bar o)}\): un completamento a base di \(\mathbb R^n\).
\end{itemize}

Poniamo quindi
\[
T=\begin{bmatrix}T_o & T_{\bar o}\end{bmatrix},
\qquad
x=Tz.
\]

\subsection{Forma del sistema trasformato}

Nelle nuove coordinate:
\[
z^+=T^{-1}AT\,z+T^{-1}Bu,
\qquad
y=CTz+Du.
\]

Poiché \(\bar{\mathcal O}\) è \(A\)-invariante, la matrice trasformata assume la struttura a blocchi
\[
\widetilde A=
\begin{bmatrix}
A_o & 0\\
A_{\bar o o} & A_{\bar o}
\end{bmatrix}.
\]
Inoltre, poiché tutte le colonne di \(T_{\bar o}\) appartengono a \(\ker(C)\), si ha
\[
CT=
\begin{bmatrix}
C_o & 0
\end{bmatrix}.
\]

Quindi il sistema assume la \textbf{forma standard di osservabilità}:
\[
\boxed{
\begin{cases}
z_o^+ = A_o z_o + B_o u,\\
z_{\bar o}^+ = A_{\bar o o} z_o + A_{\bar o} z_{\bar o} + B_{\bar o} u,\\
y = C_o z_o + Du.
\end{cases}
}
\]

In forma matriciale:
\[
\boxed{
z^+=
\begin{bmatrix}
A_o & 0\\
A_{\bar o o} & A_{\bar o}
\end{bmatrix} z
+
\begin{bmatrix}
B_o\\
B_{\bar o}
\end{bmatrix}u,
\qquad
y=
\begin{bmatrix}
C_o & 0
\end{bmatrix}z+Du.
}
\]

\subsection{Interpretazione}

La parte \(z_o\) è la parte \textbf{osservabile} del sistema.  
La parte \(z_{\bar o}\) è la parte \textbf{non osservabile}: non compare direttamente in uscita.

Anzi, poiché nella legge di uscita compare solo \(z_o\),
\[
y=C_oz_o+Du,
\]
il blocco \(z_{\bar o}\) non può essere ricostruito dalle sole misure di uscita.

Inoltre, il sottosistema \((A_o,C_o)\) è completamente osservabile.

\medskip

Gli autovalori di \(A_o\) sono gli \textbf{autovalori osservabili}, mentre quelli di \(A_{\bar o}\) sono gli \textbf{autovalori non osservabili}.

%=========================================================
\section{Parallelismo completo con la forma standard di raggiungibilità}

A questo punto il parallelismo è perfetto.

\subsection*{Forma standard di raggiungibilità}
\[
\widetilde A=
\begin{bmatrix}
A_R & A_{RN}\\
0 & A_N
\end{bmatrix},
\qquad
\widetilde B=
\begin{bmatrix}
B_R\\
0
\end{bmatrix}.
\]
Qui la parte \emph{non raggiungibile} non è direttamente influenzabile dall'ingresso.

\subsection*{Forma standard di osservabilità}
\[
\widetilde A=
\begin{bmatrix}
A_o & 0\\
A_{\bar o o} & A_{\bar o}
\end{bmatrix},
\qquad
\widetilde C=
\begin{bmatrix}
C_o & 0
\end{bmatrix}.
\]
Qui la parte \emph{non osservabile} non è direttamente visibile in uscita.

\medskip

Questa simmetria non è casuale: deriva dalla dualità tra osservabilità e raggiungibilità.

%=========================================================
\section{Conclusioni della lezione}

In questa lezione abbiamo fissato le idee fondamentali su osservabilità e ricostruibilità.

\begin{enumerate}
    \item In tempo discreto, la matrice di osservabilità in \(t\) osservazioni è
    \[
    \mathcal O_t=
    \begin{bmatrix}
    C\\
    CA\\
    \vdots\\
    CA^{t-1}
    \end{bmatrix},
    \]
    e il sottospazio di inosservabilità in \(t\) passi è
    \[
    \bar{\mathcal O}_t=\ker(\mathcal O_t).
    \]

    \item La catena
    \[
    \bar{\mathcal O}_1 \supseteq \bar{\mathcal O}_2 \supseteq \cdots
    \]
    è decrescente e si stabilizza al più dopo \(n\) passi.

    \item La ricostruibilità in \(t\) passi è equivalente alla condizione
    \[
    \ker(\mathcal O_t)\subseteq \ker(A^t).
    \]

    \item In tempo continuo, ricostruibilità e osservabilità coincidono; in tempo discreto no, salvo casi particolari, ad esempio quando \(A\) è invertibile.

    \item L'osservabilità è una proprietà strutturale, invariante per cambi di coordinate.

    \item Il sottospazio di inosservabilità
    \[
    \bar{\mathcal O}=\ker(\mathcal O)
    \]
    è il più grande sottospazio \(A\)-invariante contenuto in \(\ker(C)\).

    \item In opportune coordinate, il sistema si separa in una parte osservabile e una non osservabile:
    questa è la \textbf{forma standard di osservabilità}.
\end{enumerate}

\bigskip

La lezione successiva potrà proseguire in modo naturale verso:
\begin{itemize}
    \item gli autovalori osservabili e non osservabili;
    \item il lemma PBH per l'osservabilità;
    \item gli osservatori di stato e, in particolare, l'osservatore di Luenberger.
\end{itemize}